%% =====================================================================
%%  Mode d'emploi
%% =====================================================================
\thispagestyle{otFront}
\phantomsection\pdfbookmark[0]{Mode d'emploi}{bm-mode-emploi}
\otPageTitre{Mode d'emploi}{Comment ce livre est organisé}

%% --- Les quatre rubriques ------------------------------------------
\otIntertitre{Chaque chapitre, une synthèse puis quatre rubriques}
{\fontsize{10}{13.5}\selectfont
La synthèse \textbf{L'essentiel} ouvre chaque chapitre : définitions,
théorèmes, démonstrations à connaître, méthodes et pièges. Les
vingt-cinq exercices suivent, classés en quatre rubriques.\par}
\vspace{2.5mm}
\begin{tcbraster}[raster columns=2,raster equal height,raster column skip=4mm,raster row skip=2.5mm]
\tcbset{colback=otNavyTint,colframe=otNavyTint,boxrule=0pt,arc=2mm,
  left=3.5mm,right=3mm,top=2.5mm,bottom=2.5mm}%
\begin{tcolorbox}
  \otRond[otNavy]{1}\hspace{1.5mm}{\otSemi\fontsize{7.6}{10}\selectfont\color{otNavy}\otCaps[12]{Applications}}\par
  \vspace{1.2mm}{\raggedright\fontsize{9.5}{12.5}\selectfont Huit exercices pour vérifier le cours.}
\end{tcolorbox}
\begin{tcolorbox}
  \otRond[otNavy]{2}\hspace{1.5mm}{\otSemi\fontsize{7.6}{10}\selectfont\color{otNavy}\otCaps[12]{Entraînement}}\par
  \vspace{1.2mm}{\raggedright\fontsize{9.5}{12.5}\selectfont Dix exercices de méthode, niveau colles et DS.}
\end{tcolorbox}
\begin{tcolorbox}
  \otRond[otNavy]{3}\hspace{1.5mm}{\otSemi\fontsize{7.6}{10}\selectfont\color{otNavy}\otCaps[12]{Approfondissement}}\par
  \vspace{1.2mm}{\raggedright\fontsize{9.5}{12.5}\selectfont Six problèmes dans l'esprit des concours.}
\end{tcolorbox}
\begin{tcolorbox}[colback=otGreenTint,colframe=otGreenTint]
  \otRond[otGreenDark]{4}\hspace{1.5mm}{\otSemi\fontsize{7.6}{10}\selectfont\color{otGreenDark}\otCaps[12]{Défi}}\par
  \vspace{1.2mm}{\raggedright\fontsize{9.5}{12.5}\selectfont Un problème de synthèse, vers un résultat célèbre.}
\end{tcolorbox}
\end{tcbraster}

\vspace{4.5mm}
%% --- Anatomie d'un exercice ----------------------------------------
\otIntertitre{Anatomie d'un exercice}
\begin{tcolorbox}[colback=white,colframe=otRule,boxrule=0.4pt,arc=2mm,
  left=5mm,right=5mm,top=9mm,bottom=4mm]
\begin{tikzpicture}
  \useasboundingbox (0,-9mm) rectangle (\linewidth,4mm);
  \node[anchor=west,fill=otNavy,rounded corners=1.6pt,inner sep=0pt,minimum width=7mm,
    minimum height=6.2mm,text=white,font=\otBold\fontsize{10}{10}\selectfont] (num) at (0,0) {17};
  \node[anchor=west,inner sep=0pt,text=otInk,font=\otSemi\fontsize{10.5}{13}\selectfont]
    (tit) at (10.5mm,0) {Images directes et images réciproques};
  \node[anchor=east,inner sep=0pt,text=otMuted,font=\otMedium\fontsize{7}{9}\selectfont]
    (ref) at (\linewidth,0) {Corrigé p.~64};
  \node[anchor=east,inner sep=0pt] (dots) at ([xshift=-3mm]ref.west) {\otDots{2}};
  \node[anchor=west,inner sep=0pt,text=otInk,font=\fontsize{10.5}{13}\selectfont]
    at (10.5mm,-7mm) {Soit $f:E\to F$ une application et $A$, $B$ deux parties de $E$.};
  \foreach \n/\c in {1/num,2/tit,3/dots,4/ref}{%
    \node[circle,fill=otGreen,inner sep=0pt,minimum size=4.4mm,text=white,
      font=\otBold\fontsize{7}{8}\selectfont] at ([yshift=6.5mm]\c.north) {\n};
    \draw[otGreen,line width=0.5pt] ([yshift=4.3mm]\c.north) -- ([yshift=0.6mm]\c.north);}
\end{tikzpicture}
\end{tcolorbox}
\vspace{2mm}
{\fontsize{10}{13.5}\selectfont
\begin{tabular}{@{}p{78mm}@{\hspace{6mm}}p{78mm}@{}}
\otRond{1}\enspace Le numéro, à reporter dans le tableau de bord. &
\otRond{3}\enspace Le niveau de difficulté, de 1 à 3 points.\\[2mm]
\otRond{2}\enspace Le titre, qui annonce la notion travaillée. &
\otRond{4}\enspace La page du corrigé (lien cliquable dans le PDF).
\end{tabular}}

\vspace{4.5mm}
%% --- Les corrigés ---------------------------------------------------
\otIntertitre{Des corrigés qui expliquent}
\begin{minipage}[c]{96mm}
\begin{tcolorbox}[colback=white,colframe=otRule,boxrule=0.4pt,arc=2mm,
  left=5mm,right=5mm,top=4mm,bottom=4mm]
\otLabel{1.}\enspace\methode{raisonner par double inclusion}\\[1mm]
\hspace*{1.4em}$f^{-1}(F\setminus B)=\res{E\setminus f^{-1}(B)}$
\end{tcolorbox}
\end{minipage}\hfill
\begin{minipage}[c]{62mm}
\fontsize{10}{13.5}\selectfont
\otRond{5}\enspace La méthode à retenir, rappelée avant chaque raisonnement.\par\vspace{2.5mm}
\otRond{6}\enspace Le résultat final, surligné en vert.
\end{minipage}

\vspace{4.5mm}
%% --- Les évaluations ------------------------------------------------
\otIntertitre{Seize sujets pour s'évaluer}
{\fontsize{10}{13.5}\selectfont
La partie \textbf{Évaluations} propose seize sujets inédits — six DS, quatre
DM, deux examens et quatre sujets type concours — corrigés avec un barème
détaillé.\par}

\vspace{4.5mm}
%% --- En ligne -------------------------------------------------------
\begin{tcolorbox}[colback=otNavyTint,colframe=otNavyTint,boxrule=0pt,arc=2mm,
  left=5mm,right=5mm,top=3mm,bottom=3mm]
\begin{minipage}[c]{24mm}
{\color{otNavy}\qrcode[height=18mm,nolink]{\otURL}}
\end{minipage}\hspace{5mm}%
\begin{minipage}[c]{120mm}
{\otSemi\fontsize{11}{14}\selectfont\color{otNavy}Et aussi en ligne}\par\vspace{1.5mm}
{\fontsize{10}{13.5}\selectfont Retrouvez les cours, les fiches de synthèse et d'autres
exercices sur \textbf{\otSite}. Scannez le code avec votre téléphone pour y accéder.}
\end{minipage}
\end{tcolorbox}
\clearpage
